% ============================================================
% Заключение
% ============================================================
\chapter*{Заключение}
\phantomsection\addcontentsline{toc}{chapter}{Заключение}

В~настоящем учебном пособии систематизированы теоретические основы и~методы анализа дискретного микрорельефа рабочей поверхности деталей машин применительно к~подшипникам скольжения.

В~первой главе рассмотрены общие сведения о~подшипниках скольжения: классификация, принцип работы гидродинамического подшипника, режимы трения, конструкции и~области применения. Систематизированы виды детерминированных микроструктур рабочей поверхности, включая их~типы, формы, параметры, схемы расположения и~технологии формирования. Показано, что влияние микроструктур на~трибологические характеристики обусловлено генерацией дополнительного гидродинамического давления, удержанием смазки при~пусковых режимах и~захватом частиц износа. Рассмотрена связь поверхностного текстурирования с~концепцией зелёной трибологии.

Во~второй главе сформулирована единая математическая основа для~численного моделирования смазочного слоя. Для~трёх типовых постановок (радиальный подшипник, упорный подшипник, обратная пара) введены согласованные поверхностные координаты, масштабы безразмеризации и~унифицированная форма уравнения Рейнольдса. Сформулирована параметрическая модель микрорельефа и~массосохраняющая постановка кавитации (JFO).

В~третьей главе описана численная реализация решения задач гидродинамики смазочного слоя: консервативная дискретизация, метод Red-Black SOR на~GPU, алгоритм решения задачи кавитации JFO с~двухуровневой итерационной структурой.

В~четвёртой главе проанализировано влияние дополнительных эффектов на~характеристики смазочного слоя: нелинейность (пьезовязкость и~сжимаемость), термовязкостные эффекты и~шероховатость поверхности. Сформулирована иерархия моделей-расширений базовой изотермической постановки.

В~пятой главе приведены три~подробных примера инженерного расчёта подшипников скольжения с~микрорельефом: радиальный подшипник центробежного насоса, радиальный подшипник шарошечного долота и~упорный подшипник с~секторными колодками. Каждый пример демонстрирует полный цикл расчёта~-- от~постановки задачи и~выбора параметров микроструктур до~анализа и~интерпретации результатов. Показано, что оптимально подобранные микроструктуры способны существенно повысить несущую способность подшипника и~снизить потери на~трение.

\textbf{Практическое значение.} Полученные знания и~навыки могут быть использованы при~проектировании и~модернизации подшипниковых узлов технологических машин и~оборудования нефтегазовой, энергетической и~других отраслей. Методика расчёта, освоенная при~выполнении курсовой работы, позволяет обучающемуся самостоятельно оценивать эффективность применения микрорельефа для~конкретных условий эксплуатации.

\textbf{Перспективы развития.} Направление анализа и~оптимизации микрорельефа рабочих поверхностей продолжает активно развиваться. К~перспективным задачам относятся: учёт термогидродинамических эффектов в~расчётах текстурированных подшипников, многокритериальная оптимизация параметров микроструктур с~применением методов машинного обучения, разработка новых технологий формирования микрорельефа (аддитивные технологии, ультракороткоимпульсная лазерная обработка), а~также расширение области применения текстурирования на~другие узлы трения~-- уплотнения, направляющие, зубчатые передачи.
